\chapter{Who would use it, and why}
\label{ch:adoption-scale}

An organization may like the demonstration and still decline the product. Its
authors have obligations, its security team has a privacy boundary, and its
readers already have a review culture. The commercial question after the
feasibility case is therefore concrete: who will use Humanvoice, what will it
replace, and what evidence would justify moving from one document to a
portfolio?

The repository does not yet contain a defensible market-size estimate. It does
contain a clear first job. A high-consequence document has a named expert
reader, several expensive review rounds, and equations, citations, or other
objects that a general editor cannot safely handle. Adoption should begin where
that job is already visible and measurable.

Figure~\ref{fig:value-mechanism} gives the commercial proposition in one line:
saved attention creates value only after meaning and reader burden pass the
guard.

\hvFigureValueMechanism

\section{The buyer, the author, and the reader}
\label{sec:adoption-roles}

The buyer may sign the contract, the author may do the editing, and the reader
may decide whether the document matters. Treating those three people as one
``user'' would optimize the wrong outcome.

\begin{table}[H]
\centering
\small
\begin{tabularx}{\textwidth}{@{}p{0.18\textwidth}p{0.27\textwidth}L@{}}
\toprule
Role & Immediate job & Evidence that adoption is real \\
\midrule
Sponsor or buyer & Fund a bounded workflow and pay for reduced delay,
lower incident risk, or a better decision record. & Approves the comparator,
privacy boundary, and continuation threshold. \\
Author or analyst & Revise a source without losing equations, numbers,
citations, or ownership of the claim. & Opens findings, records decisions, and
returns to the incumbent path when the system abstains. \\
Editor or methods lead & Make a document's argument teachable and its review
repeatable across authors. & Reuses contracts, fixtures, and reader rubrics
without turning them into a style score. \\
Expert reader or referee & Reconstruct the problem, mechanism, evidence,
uncertainty, and requested action under a real time budget. & Gives a located
decision, confidence, and change trigger. \\
Security or policy owner & Decide what source and response data may leave the
institution and what disclosure is required. & Signs the boundary and can
replay or revoke a remote call. \\
\bottomrule
\end{tabularx}
\caption{Different roles make different adoption claims.}
\label{tab:adoption-roles}
\end{table}

The buyer's outcome is not ``more human-sounding text.'' It is a document that
can be accepted, challenged, or funded on the reader's own terms. The author's
outcome is not a lower diagnostic count. It is less time spent rediscovering
why a reader objected, with enough provenance to decide whether the objection
was substantive. The reader's outcome is not agreement. It is a well-founded
decision with visible uncertainty.

\begin{cardexample}{A buyer can reject a successful demo}
An engineering team demonstrates that a model can turn a paragraph into fluent
prose. The sponsor asks whether the revision preserves a denominator, a
qualification, and the identity of a cited study. The team has no source map,
no protected comparison, and no way to show the reader's time. The demo may be
linguistically impressive, but it has not addressed the buyer's job.
\end{cardexample}

The invalid inference is that a positive author reaction is sufficient for
adoption. The first pilot must observe all three roles where possible. If a
document has no consequential reader, it is a useful parser fixture but a
weak product case.

\section{Three beachheads, chosen by the job}
\label{sec:adoption-beachheads}

Start with settings in which the job is visible, not with an abstract total
addressable market. Three settings fit the current evidence and differ enough
to test whether the mechanism travels.

\subsection{Research and methods groups}

A research group produces a methods note, working paper, or monograph chapter
whose reader is a technical referee or collaborator. The source may contain
custom mathematics, code, citations, and a long chain of assumptions. The
group already uses version control and a build system, so the first product
integration can be local and command-line driven. The value is a shorter path
from a review comment to a verified revision, not automatic copy-editing.

The risk is that a researcher may treat the tool as an additional referee and
accept a plausible but technically wrong suggestion. The pilot therefore
starts with deterministic source correspondence and a human-controlled
finding queue. A model adapter is optional and must be disclosed.

\subsection{Investment, policy, and program offices}

An investment or policy office produces a decision document for an executive
who is expert but time-constrained. The document may mix quantitative results,
institutional constraints, and a request for funding or a decision. Here
the reading brief and unfamiliar-reader exercise are especially valuable: the
question is whether a decision-maker can reconstruct the ask without private
meetings.

The risk is confidentiality. A remote model call, a retained page image, or a
shared benchmark fixture can expose information that is more sensitive than
the prose itself. This setting is therefore a strong test of the local-first
boundary and of the continuity requirement: the office must be able to finish
the document if Humanvoice is unavailable.

\subsection{Publishers, funders, and research-support units}

A publisher or funder may not write the document, but it can provide a
reading-brief and review service across many authors. The buyer's value is
portfolio consistency and a reusable evidence record; the author's value is a
clearer review path; the reader's value is a document that can be evaluated
without learning the organization's internal workflow.

The risk is scale before validity. A support unit can process many documents
and still measure only surface cleanliness. This setting should be a second
stage, after the single-document feasibility run has fixed the object model,
reader rubric, and disclosure language.

\begin{table}[H]
\centering
\small
\begin{tabularx}{\textwidth}{@{}p{0.20\textwidth}p{0.24\textwidth}p{0.24\textwidth}L@{}}
\toprule
Setting & Strong reason to start & First measurable outcome & Main stop signal \\
\midrule
Research group & Protected technical objects and existing source control. &
Time from located finding to accepted revision; protected-object incidents. &
Authors cannot distinguish a useful question from a model rewrite. \\
Investment or policy office & High review cost and a named executive decision. &
Reconstruction by a reader unfamiliar with the project, followed by feedback rounds to a decision. &
Privacy or disclosure boundary cannot be signed. \\
Publisher or support unit & Repeated documents and a reusable reader rubric. &
Portfolio-level burden and reader consistency after single-document validity. &
Scale creates a surface score with no reader endpoint. \\
\bottomrule
\end{tabularx}
\caption{Beachheads and the evidence each one can produce.}
\label{tab:beachheads}
\end{table}

The proposal should name one live setting for the twelve-week build. It may
collect fixtures from the other two, but it should not promise three distinct
markets with one underpowered pilot. A focused first customer gives the
comparator, privacy rules, and reader role enough specificity to be measured.

\section{The alternatives Humanvoice must beat}
\label{sec:adoption-alternatives}

If Humanvoice is absent, the work does not disappear. An author, domain editor,
general language model or style tool, version control, a build system, citation
checks, and one or more expert readers form the incumbent bundle. Humanvoice
must earn its place by improving the handoff between those activities.

\begin{table}[H]
\centering
\small
\begin{tabularx}{\textwidth}{@{}p{0.22\textwidth}p{0.25\textwidth}p{0.25\textwidth}L@{}}
\toprule
Alternative & What it does well & What remains unconnected & Humanvoice's
wedge \\
Manual expert editing & Interprets claims and audience in context. & Expensive,
hard to replay, and often loses the reason for a revision. & Preserve the
finding, source span, protected objects, and reader response around the human
judgment. \\
General language model & Generates fluent alternatives and questions quickly.
& May widen claims, obscure ownership, or require unapproved remote processing.
& Use bounded questions and explicit acceptance rather than silent rewriting. \\
CI linter or style guide & Replays deterministic checks at low cost. & Usually
cannot connect a warning to the reader's decision or technical object. & Keep
narrow diagnostics and attach them to a reading brief. \\
Version control and diff & Shows that source changed and permits rollback. & A
character-level diff does not say whether a denominator, citation, or claim
scope changed. & Add object-aware comparison and a reason for each revision. \\
External review service & Supplies a reader and an acceptance signal. &
Feedback may be unstructured and source provenance may be lost. & Return a
replayable packet that links reader action to author work. \\
\bottomrule
\end{tabularx}
\caption{The product is a connective layer over an incumbent bundle.}
\label{tab:adoption-alternatives}
\end{table}

The wedge is strongest where the alternatives are each useful but their
interfaces are weak. It is not a claim that Humanvoice is a better editor,
better model, or better referee. It is a claim that a named reader's decision
can be connected to a protected revision record owned by the author.

\section{A staged adoption path}
\label{sec:adoption-path}

Adoption should proceed in the same order as the evidence ladder. Each stage
has a customer-visible deliverable and a reason to stop. This makes expansion
a sequence of earned options rather than a commitment to a large platform
before the first object map works.

\begin{enumerate}
\item \textbf{Instrument one document.} Agree on the reading brief,
      incumbent bundle, privacy class, protected-object set, and acceptance
      endpoint. Produce a baseline run even if no Humanvoice finding is shown.
\item \textbf{Run the vertical slice.} Parse, locate findings, record author
      decisions, compare protected objects, and prepare a reader packet for
      someone without project history. Keep the model adapter off unless the owner explicitly permits
      it.
\item \textbf{Repeat within one setting.} Add documents from the same genre
      and reader role. Estimate burden, abstention, and heterogeneity before
      comparing settings. A second document is a replication opportunity, not
      proof of generality.
\item \textbf{Compare the incumbent bundle.} Counterbalance or randomize
      preparation and reading order where practical. Declare the primary
      endpoint and preserve protocol deviations.
\item \textbf{Offer a portfolio service.} Only after the comparative result
      supports the workflow should a support unit reuse contracts, fixtures,
      and reader instruments across authors.
\end{enumerate}

At each stage, the customer should receive a short decision packet: what was
run, what changed, what was protected, what the reader did, how much time it
took, and what remains uncertain. A dashboard that reports only a finding
count is not a product milestone. The customer must be able to decide whether
to continue using the incumbent, add a human editor, or fund the next stage.

\paragraph{The adoption handoff.}
Before moving from one document to a portfolio, ask a person who did not build
the fixtures to run the packet. If they cannot identify the reading brief,
replay the protected comparison, and explain an abstention, the product is
not ready for scale. More documents would multiply an unresolved interface
failure.

\section{Pricing as a testable hypothesis}
\label{sec:adoption-pricing}

Pricing should follow the unit of value observed in the pilot. A per-token or
per-page price would reward volume even when the system creates review work.
The first offer should therefore be a bounded service or pilot priced around a
document cohort, an agreed reader exercise, and a retained evidence packet.
Only after repeat use is measured should a recurring subscription or
institutional license be considered.

Let $n$ be the number of documents in a cohort, $s$ the service and support
cost per document, $K_0$ the fixed onboarding cost, and $W$ the observed value
of avoided or improved review work. A pilot contribution can be written as
\[
  \Pi = n(W-s)-K_0.
\]
This is a planning identity, not a forecast. $W$ must be estimated from the
reader and author logs in \autoref{app:cost-worksheets}; it should include an
incident-cost range when a missed qualification or protected-object change is
material. A customer should not be charged for a benefit the pilot never
measured.

\begin{table}[H]
\centering
\small
\begin{tabularx}{\textwidth}{@{}p{0.24\textwidth}p{0.27\textwidth}L@{}}
\toprule
Offer form & What the customer receives & Why it is appropriate at this stage \\
\midrule
Feasibility service & One live run, protected comparison, reader exercise,
and decision packet. & Buys learning while the object and reading briefs are
still being tested. \\
Pilot cohort & A fixed number of same-setting documents, fixtures, and a
comparative report. & Measures repeatability and burden without claiming a
general platform. \\
Institutional deployment & Local service, support, retention controls, and
portfolio rubric. & Requires evidence that the workflow works across authors
and that operations are acceptable. \\
\bottomrule
\end{tabularx}
\caption{Pricing forms should match the evidence stage.}
\label{tab:pricing-forms}
\end{table}

The proposal should resist a large subscription forecast until there is
observed renewal behavior. A financier can still evaluate the option value of
the MVP: a modest first investment buys a working source map, a corpus, and a
credible estimate of $W$. If those objects do not materialize, the project
has learned early and can stop without having built an expensive service
whose value is only asserted.

\section{Operating model and trust boundary}
\label{sec:adoption-operating}

An adopted service needs an owner for every handoff. The author owns the
substantive revision. The editor or methods lead owns the reading brief and
rubric. The engineer owns parser and build failures. The security or policy
owner owns the processing boundary. The sponsor owns the continuation rule.
These roles may be held by two people in a small pilot, but the decisions must
remain distinct in the records.

The local-first path should be the default operating mode. Deterministic
parsing, source maps, protected comparisons, and narrow diagnostics run in
the repository boundary. A model adapter, if allowed, receives only the
smallest excerpt needed for a declared question and returns a finding rather
than a replacement document. The run records provider, version, prompt or
configuration identifier, retention, and abstention. A customer can then
choose a fully local path without losing the core product.

Continuity protects the customer's deadline. If a
remote provider is unavailable, the author should be able to continue editing
the source, build the document, and hand the reader the incumbent packet. The
eventual Humanvoice run can import the parent snapshot and show what changed.
This protects the customer's deadline and makes the product's value easier to
trust: adoption does not require surrendering control of the document.

\begin{table}[H]
\centering
\small
\begin{tabularx}{\textwidth}{@{}p{0.25\textwidth}p{0.30\textwidth}L@{}}
\toprule
Operating promise & Minimum customer-visible behavior & Failure response \\
\midrule
Local processing & Source, pages, and deterministic findings stay inside the
declared boundary. & Block or quarantine an unapproved remote call. \\
Human ownership & No revision enters the reader packet without an author
decision and protected comparison. & Mark the run incomplete and preserve the
parent source. \\
Visible uncertainty & Unsupported objects and model disagreement appear as
abstentions. & Route to the incumbent workflow or a named expert. \\
Continuity & The document can be completed if a component is unavailable. &
Export the source and reading brief; import the later decision. \\
Accountable records & A sponsor can inspect the evidence behind a continuation
or stop decision. & Retain the minimum run, response, and rationale records. \\
\bottomrule
\end{tabularx}
\caption{Operating promises that make adoption reversible.}
\label{tab:operating-promises}
\end{table}

\section{The scale decision}
\label{sec:adoption-scale-decision}

At the end of the first comparative cohort, the sponsor should make a scoped
scale decision, not a binary verdict on ``AI writing.'' Four results are
possible:

\begin{description}
\item[Scale within the setting.] The reader endpoint improves or the burden
  falls, protected-object incidents remain controlled, and the privacy and
  continuity promises hold for the tested population.
\item[Deepen the instrument.] The workflow is usable but the endpoint is too
  noisy, the corpus is too narrow, or the reading brief is underspecified.
  Invest in measurement before adding features.
\item[Reposition.] Deterministic source correspondence or comparison is
  valuable, but the complete packet does not change reader outcomes. Offer an
  internal editing and evidence service without claiming efficacy.
\item[Stop.] The product creates burden, cannot protect the source objects,
  violates the privacy boundary, or cannot produce a credible comparison.
\end{description}

The scale rule should be written before the result is seen. For example, move
to a second setting only if the held-out documents meet the protected-object
threshold, the median author burden is no higher than the declared margin,
and the reader can state the action and uncertainty above a predeclared
threshold. The exact values belong to the study protocol; the important point
is that a fluent demo cannot waive them.

\begin{plainbox}{What the proposal is buying}
The requested investment buys an option, not a promise of a market. It buys a
working path from source snapshot to reader decision, a corpus that can expose
where the path fails, and an economic worksheet tied to observed review work.
Those objects make the next financing decision more informed. They also make
it possible to stop without confusing a polished prototype with a validated
product.
\end{plainbox}

The next appendices specify the records, fixtures, reader forms, and cost
inputs that make this scale rule executable. They are part of the product
proposal because an operating service cannot be separated from the evidence
that tells a sponsor whether to keep paying for it.
